\section{Boundary multiplicity obstructs the actual character source}
\label{sec:character-boundary-obstruction}

The strict-exterior result MA-MO-024 leaves open what happens
on the edge itself. Integrating its actual translated source
orbit gives a stronger estimate than testing one translation
at a time. At the equality threshold the decisive convergence
is weak convergence in the original Hilbert norm.

Retain the nontrivial character, structural place set, actual
Fourier pair, real cutoff, and original smooth source action
of (MA-MO.550)--(MA-MO.552). In particular
\[
 \begin{gathered}
 H=H_{\Sigma,a,s}=L^2(\mathbb R,w_{a,s}(x)\,dx),
       \qquad 0<a<1/2,\quad s\ge0,\\
 w_{a,s}(x)=(e^{2ax}+e^{-2ax})(1+x^2)^{s/2},\qquad
 U_th(x)=h(x-t),\\
 D_g^{\,0}=\mathscr D^{\,\rm raw}_{\eta,g}:
                         \mathcal E\longrightarrow\mathcal E .
 \end{gathered}
 \tag{MA-MO.600}
\]
The last map is defined on the actual smooth common sector
by the original wide inverse and the original source Fourier
maps. No bounded narrow extension is assumed in its definition.

\paragraph{MA-MO-025 (the exact deficient-boundary obstruction).}
Let \(\rho\) be a completed nontrivial zero of \(L_K(s,\eta)\),
of multiplicity \(m\), such that
\[
             |\Re\rho-1/2|=a,\qquad s\le 2m-1 .
 \tag{MA-MO.601}
\]
For every nonzero complex \(g\in C_c^\infty(\mathbb R)\),
the actual map \(D_g^{\,0}\) has no bounded extension
\(H\to H\). The conclusion includes the equality
\(s=2m-1\), and has no exception when the Mellin transform
of \(g\) vanishes at the zero. It also holds within the
nonzero test class \(g(0)=\mathcal Mg(\pm1/2)=0\).

Consequently, boundedness of this actual full source action
for even one nonzero compact test requires full retention
of every nontrivial zero, with its complete multiplicity,
in the weighted partial-strip quotient. This is a necessary
condition. No fixed-width sufficiency without an interior
margin is asserted.

\begin{proof}
\emph{A boundary coefficient that survives the smoothing.}
First suppose \(\lambda=\rho-1/2=a+i\gamma\).
Use the notation \(A_\eta,H_h,V_g,B_v,\mathcal R_{h,v}\)
from (MA-MO.551) and (MA-MO.558)--(MA-MO.567).
The nonannihilation proof (MA-MO.568)--(MA-MO.572)
applies at this boundary point as well: none of its steps
requires the strict inequality \(\Re\lambda>a\).
For clarity, the function used there is
\[
 b_\lambda(x)=e^{-\lambda x}
     \sum_{\mathfrak a\ {\rm prime\ to}\ S}
       c_{\eta,\mathfrak a}(N\mathfrak a)^{-1/2-\lambda}
                              \phi(x+\log N\mathfrak a).
 \tag{MA-MO.602}
\]
It vanishes to the right of the last support edge \(r\)
of \(\phi\). In the interval \(r-\log2<x<r\), the
unit ideal alone contributes, so
\(b_\lambda=e^{-\lambda x}\phi(x)\) is not identically
zero. For any
\(\sigma_0>\max(1,1/2+\Re\lambda)\), the finite ideal
sum gives the left bound
\(\lvert b_\lambda(x)\rvert\le
C_\lambda e^{-(\sigma_0-1/2)x}\).
Its one-sided Mellin transform therefore exists on a
nonempty right half-plane. If \(C_gb_\lambda=0\), the
product of that transform and the nonzero entire function
\(\mathcal Mg\) would vanish there. Holomorphic uniqueness
and then Fourier uniqueness would give \(b_\lambda=0\),
a contradiction. This argument permits arbitrary complex
\(g\ne0\).

Choose \(v\in C_c^\infty\) detecting \(C_gb_\lambda\),
and choose \(h=e^{-\lambda x}h_0(x)\), where
\(0\ne h_0\in C_c^\infty\) is nonnegative. The bilinear
transpose \(v_g=C_{g^\vee}v\), without complex conjugation
of \(g\), gives
\[
 B_v(\lambda)=\int v(x)C_gb_\lambda(x)\,dx\ne0,\qquad
 H_h(\lambda)=\int h_0(x)\,dx>0 .
 \tag{MA-MO.603}
\]
Fix \(\operatorname{supp}h\subset[\alpha,\beta]\).
Increase \(t_0>0\) so that all support conditions
(MA-MO.561) hold and \(\beta+t_0\ge1\).
The actual dual source term then vanishes identically
for \(U_th\), \(t\ge t_0\), before its Fourier map is
applied. With the continuous output functional
\(\ell_v(y)=\int y(x)v(x)\,dx\), put
\[
 f(t)=\ell_v(D_g^{\,0}U_th),\qquad
 \mathfrak F(w)=\int_{t_0}^{\infty}e^{-wt}f(t)\,dt
                                      \quad(\Re w>1/2).
 \tag{MA-MO.604}
\]
The original source sums, the finite theta pairing, and
absolute wide Euler convergence give (MA-MO.567).
In particular its remainder \(\mathcal R_{h,v}\) and
both numerator factors in (MA-MO.603) are entire.

\emph{What a bounded extension would imply.}
Suppose that a bounded \(D:H\to H\) agrees with
\(D_g^{\,0}\) on \(\mathcal E\).
The exact weighted translation bound is
\[
 \|U_th\|_H\le
       2^{s/4}(1+t^2)^{s/4}e^{at}\|h\|_H
                                             \qquad(t\ge0).
 \tag{MA-MO.605}
\]
It follows that \(f(t)=O(e^{at}(1+t)^{s/2})\), and
dominated differentiation makes the actual integral
\(\mathfrak F\) holomorphic on \(\Re w>a\).
The cleared identity (MA-MO.574), initially proved by
wide Euler sums, extends by holomorphic uniqueness:
\[
 A_\eta(w)\bigl[
 H_h(w)V_g(-w)-\mathcal R_{h,v}(w)-\mathfrak F(w)
                 \bigr]=H_h(w)B_v(w),\qquad \Re w>a .
 \tag{MA-MO.606}
\]
There is no use here of a reciprocal multiplier on the
boundary line.

The finite Euler deletions have no zeros when
\(\Re(1/2+w)>0\); the gamma factors are finite and
nonzero at the present completed zero. Thus \(A_\eta\)
has exactly the same multiplicity \(m\) at \(\lambda\).
Write
\[
 A_\eta(w)=a_m(w-\lambda)^m+
                      O((w-\lambda)^{m+1}),\qquad
 a_m=\frac{A_\eta^{(m)}(\lambda)}{m!}\ne0 .
 \tag{MA-MO.607}
\]
In a sufficiently small disk \(\lambda\) is its only
zero. Evaluating (MA-MO.606) at
\(w_\delta=\lambda+\delta\), \(0<\delta<\delta_0\),
and using the nonzero numerator in (MA-MO.603), gives
\[
 \lim_{\delta\downarrow0}
       \delta^m\mathfrak F(\lambda+\delta)
       =-\frac{H_h(\lambda)B_v(\lambda)}{a_m}\ne0 .
 \tag{MA-MO.608}
\]
The other two terms in the meromorphic expression
(MA-MO.567) are holomorphic near \(\lambda\), so they
do not affect this limit. This retains the full
multiplicity of the zero.

\emph{An actual Hilbert orbit integral and its original-core
approximants.}
For each fixed \(\delta>0\), define
\[
 \begin{aligned}
 u_{\delta,T}
     &=\int_{t_0}^{T}e^{-w_\delta t}U_th\,dt,
                         &&T>t_0,\\
 u_\delta
     &=\int_{t_0}^{\infty}e^{-w_\delta t}U_th\,dt .
 \end{aligned}
 \tag{MA-MO.609}
\]
These are Bochner integrals in \(H\).
Translations are strongly continuous in this exact norm:
first verify this on compact smooth functions, and then
use their density and the locally uniform operator bound
(MA-MO.605). The infinite integral converges because
its norm integrand is bounded by
\(C e^{-\delta t}(1+t)^{s/2}\).
For finite \(T\), differentiation under the compact
parameter integral shows
\[
 u_{\delta,T}\in C_c^\infty
                 ([\alpha+t_0,\beta+T])\subset\mathcal E,
 \qquad
 u_{\delta,T}\longrightarrow u_\delta\quad\hbox{in }H .
 \tag{MA-MO.610}
\]
Thus the approximating inputs lie in the original
source domain, with no inverse continuation.
Since \(D\) is assumed bounded, it commutes with these
Bochner integrals and agrees with the actual source map
on every integrand and every finite approximant.
Applying the fixed continuous functional yields
\[
 \begin{aligned}
 \ell_v(Du_{\delta,T})
       &=\int_{t_0}^{T}e^{-w_\delta t}f(t)\,dt,\\
 \ell_v(Du_\delta)&=\mathfrak F(w_\delta).
 \end{aligned}
 \tag{MA-MO.611}
\]
We do not assert that the uncontinued raw source formula
is independently defined on \(u_\delta\). Only the
hypothetical bounded extension is applied to that limit.

\emph{The exact tail in the inherited sum norm.}
Changing variables in the orbit integral gives
\[
 u_\delta(x)=e^{-w_\delta x}
       \int_{-\infty}^{\,x-t_0}e^{w_\delta y}h(y)\,dy .
 \tag{MA-MO.612}
\]
This pointwise identity also follows from the finite
integrals and their local uniform convergence.
In particular
\[
 \begin{cases}
 u_\delta(x)=0,&x\le\alpha+t_0,\\
 u_\delta(x)=H_h(w_\delta)e^{-w_\delta x},
                                      &x\ge\beta+t_0 .
 \end{cases}
 \tag{MA-MO.613}
\]
On the fixed interval between these two endpoints,
\(u_\delta\) is uniformly bounded for
\(0<\delta\le\delta_0\); the same is true on every
fixed compact set. Put \(R_0=\beta+t_0\ge1\).
Since \(H_h(w_\delta)\) is bounded for small \(\delta\),
the exact right-tail weight gives
\[
 \begin{aligned}
 \|u_\delta\|_H^2
 &\le C+
 C\int_{R_0}^{\infty}e^{-2\delta x}
             (1+e^{-4ax})(1+x^2)^{s/2}\,dx\\
 &\le C_s\delta^{-(s+1)},\qquad 0<\delta\le\delta_0\le1 .
 \end{aligned}
 \tag{MA-MO.614}
\]
For the last bound, use
\((1+x^2)^{s/2}\le 2^{s/2}x^s\) for \(x\ge1\),
and integrate \(e^{-2\delta x}x^s\) over the positive
half-line. This computation retains both exponential
terms of the sum norm.

If \(s<2m-1\), (MA-MO.614) implies
\[
       \|\delta^m u_\delta\|_H
          \le C\delta^{\,m-(s+1)/2}\longrightarrow0 .
 \tag{MA-MO.615}
\]
Boundedness of \(D\) and continuity of \(\ell_v\) force
\(\delta^m\mathfrak F(w_\delta)\to0\) by
(MA-MO.611), contradicting (MA-MO.608).

\emph{The equality threshold.}
If \(s=2m-1\), the vectors
\(z_\delta=\delta^m u_\delta\) are uniformly bounded
in \(H\), by (MA-MO.614).
They tend uniformly to zero on each fixed compact
interval, by (MA-MO.612)--(MA-MO.613), since \(m\ge1\).
For a compactly supported \(k\in H\), weighted
Cauchy--Schwarz on its support therefore gives
\(\langle z_\delta,k\rangle_H\to0\).
For arbitrary \(k\in H\), approximate \(k\) in \(H\)
by compactly supported vectors and use the uniform
bound on \(\|z_\delta\|_H\). Consequently
\[
                 \delta^m u_\delta
                       \rightharpoonup0\quad\hbox{in }H .
 \tag{MA-MO.616}
\]
The fixed functional \(\ell_v\circ D\) is continuous,
so its values on these vectors must tend to zero.
Equations (MA-MO.608) and (MA-MO.611) contradict this.
The equality case does not rely on norm convergence.

\emph{The other boundary.}
If \(\Re(\rho-1/2)=-a\), the functional equation gives
a zero of the same multiplicity for the dual character
at \(1-\rho\), whose centered real part is \(a\).
The exact source identity (MA-MO.576) is
\[
       R D_{\eta,g}^{\,0}R=D_{\bar\eta,g^\vee}^{\,0}.
 \tag{MA-MO.617}
\]
It retains the actual finite Fourier phases, and \(R\)
is an isometry of the norm (MA-MO.600).
A bounded extension on the left would therefore give
one for the right-boundary problem already contradicted.
Since \(g^\vee\ne0\), this proves the full statement.
The proof required only \(g\ne0\), so the stated moment
conditions on the test do not change the conclusion.
\end{proof}

\paragraph{The exact necessary divisor condition.}
In the centered sum norm, MA-WP-001 gives retained
multiplicity
\[
 \begin{aligned}
 b_s&=\#\{j\in\mathbb Z_{\ge0}:2j<s-1\},\\
 r_\rho&=
 \begin{cases}
 m_\rho,&|\Re\rho-1/2|<a,\\
 \min(m_\rho,b_s),&|\Re\rho-1/2|=a,\\
 0,&|\Re\rho-1/2|>a .
 \end{cases}
 \end{aligned}
 \tag{MA-MO.618}
\]
At an edge, full retention is equivalent to
\(2(m_\rho-1)<s-1\), or \(s>2m_\rho-1\).
Combining MA-MO-024 and MA-MO-025 therefore gives
\[
 \begin{gathered}
 \exists\,0\ne g\in C_c^\infty:
 D_g^{\,0}\text{ has a bounded extension on }H_{\Sigma,a,s}\\
 \Longrightarrow\quad
 r_\rho=m_\rho\quad\text{for every nontrivial completed zero}.
 \end{gathered}
 \tag{MA-MO.619}
\]
For \(0\le s\le1\), this excludes every boundary zero.
For general \(s\), a zero at the boundary is excluded
exactly when its local algebra is truncated by this
Hilbert norm. The word \emph{exactly} here describes the
proved obstruction threshold, not a sufficiency theorem.
Even when every \(r_\rho=m_\rho\), no boundedness or
trace-class estimate for the original source action is
deduced from (MA-MO.619).
The sufficient hypothesis in MA-MO-022 remains the
explicit zero-free half-plane with a strict margin.
The existing quotient traces and quotient-return maps
do not by themselves supply that missing full-source
operator estimate.

\paragraph{Source scope and finite-length placement.}
All arithmetic source identities are inherited from
the completely written MA-MO-024 proof: literal wide
ideal sums, actual structural-\(S\) source inclusion,
the original \(p_+\) map, and nontrivial-sector Poisson
summation. Their exact primary locators are Meyer
(2005), supplied \texttt{main.tex} lines 3592--3707,
2630--2675, and 3808--4043, through
(MA-MO.551)--(MA-MO.567) and (MA-MO.576).
The new argument is the actual Hilbert orbit integral,
its compact original-core approximants, and the
strong/weak limiting contradiction.
No additional zero-distribution theorem, nonvanishing
claim, or arithmetic occurrence of a boundary zero is
used.

In the inherited cote 1 trace framework, scans 079--080
and 111--112 through BR-TR-01/02, the local divisor has
finite length \(m_\rho\), while the completed quotient
retains length \(r_\rho\). These are classes in the
category of finite-length \(\mathbb C[z]\)-modules;
their underlying finite-dimensional vector spaces
carry the determinant constructions already used in
MA-WP. Equality of those lengths is necessary here
for a bounded realization of the original full source
action. It is not an analytic extension theorem, a
positivity statement, or a consequence of a motivic
weight. The newly separated boundary question is
decided at this stated threshold; the arithmetic
deficiency condition of FT002 is not established anew.
